\documentclass[11pt]{article}
\usepackage[utf8]{inputenc}
\usepackage[T1]{fontenc}
\usepackage{amsmath,amssymb,amsfonts}
\usepackage{booktabs}
\usepackage{graphicx}
\usepackage{hyperref}
\usepackage{xcolor}
\usepackage[margin=1in]{geometry}
\usepackage{multirow}
\usepackage{tabularx}
\usepackage{natbib}

\title{The Redundancy-Stability Tradeoff: Why CKA-Aware Pruning Helps SFT but Breaks RLVR in Reasoning Models}

\author{Anonymous}
\date{}

\begin{document}
\maketitle

%% ============================================================
%% ABSTRACT
%% ============================================================
\begin{abstract}
The standard pipeline for compressing reasoning LLMs---structural pruning followed by supervised fine-tuning (SFT) and reinforcement learning with verifiable rewards (RLVR)---assumes alignment is independent of pruning decisions.
We show this assumption fails: GRPO alignment on structurally pruned models produces reward improvement without accuracy gain, followed by complete training collapse.
We introduce \textbf{RASP} (Recovery-Aware Structural Pruning), a CKA-based scoring function $S(h)=I(h)\cdot(1{-}F(h))^\alpha$ that preferentially removes functionally redundant structures.
RASP achieves the best SFT recovery (61.5\% vs.\ 59.3\% for importance-only pruning on GSM8K at 20\% sparsity, $+5.75$pp at 25\%, $p{=}0.008$), yet accelerates GRPO collapse by 6$\times$ (step 10 vs.\ step 60).
This \emph{Redundancy-Stability Tradeoff}---structures safely removable for SFT provide critical exploration stability for RLVR---is replicated across two models (DeepSeek-R1-Distill-Qwen-7B, Qwen2.5-Math-7B), two sparsity levels (20\%, 25\%), and five pruning methods.
Per-structure regression confirms the mechanism: CKA redundancy $F(h)$ anticorrelates with SFT activation change ($r{=}{-}0.42$, $p{<}0.001$) but positively predicts GRPO collapse speed ($r{=}0.55$, $p{<}0.001$).
We conclude that optimal pruning is fundamentally conditioned on the recovery paradigm, and propose collapse dynamics as a zero-cost diagnostic for structural damage in pruned reasoning models.
\end{abstract}

%% ============================================================
%% 1  INTRODUCTION
%% ============================================================
\section{Introduction}

Reasoning LLMs such as DeepSeek-R1 and Qwen2.5-Math have achieved strong performance on mathematical reasoning benchmarks, but their deployment is bottlenecked by model size.
Structural pruning---removing entire attention heads or feed-forward neurons---offers wall-clock speedup without custom sparse kernels, making it a practical compression path.
The standard pipeline applies (1) pruning, (2) SFT recovery, and (3) RLVR alignment (typically GRPO) to restore and enhance reasoning ability.

Recent work claims this pipeline works well.
Self-Healing~\citep{selfhealing2025} reports that RLVR recovers pruned reasoning models to 95\%+ accuracy at 25\% compression.
GRASPrune~\citep{graspune2026} provides state-of-the-art budgeted structural pruning but does not study recovery or alignment.
No prior work has documented GRPO training collapse on pruned models, explained its mechanism, or shown that pruning strategy determines collapse dynamics.

We present the first systematic study of how structural pruning interacts with RLVR alignment in reasoning LLMs.
Our contributions are:

\begin{enumerate}
    \item \textbf{GRPO reward-accuracy decoupling is universal; pruning adds training collapse.}
    GRPO does not improve reasoning accuracy for any model post-SFT---unpruned accuracy is 81.27\% before and after GRPO ($\Delta{=}0.00$pp, identical 1072/1319 correct)---while reward increases from 0 to 0.661. Pruning introduces a distinct failure: \emph{training collapse}, where RASP-pruned models collapse permanently at step ${\sim}10$ and Standard-pruned at step ${\sim}60$, while unpruned models train stably for 100 steps.

    \item \textbf{Collapse severity diagnoses structural damage.}
    Across five pruning methods, collapse speed correlates with the degree of redundancy removal. This makes collapse dynamics a zero-cost diagnostic: fast collapse signals aggressive redundancy removal and poor RL trainability.

    \item \textbf{Recovery paradigm conditions optimal pruning.}
    RASP achieves the best SFT recovery ($+5.75$pp at 25\% sparsity, $p{=}0.008$) but the fastest GRPO collapse (6$\times$ faster than Standard). The structures optimal to prune for SFT are exactly those needed for RLVR stability---a fundamental \emph{Redundancy-Stability Tradeoff}.
\end{enumerate}

%% ============================================================
%% 2  RELATED WORK
%% ============================================================
\section{Related Work}

\paragraph{Structural pruning for LLMs.}
LLM-Pruner~\citep{llmpruner2023} and ShortGPT~\citep{shortgpt2024} remove attention heads and layers based on importance scores.
GRASPrune~\citep{graspune2026} uses gradient-based budget allocation for structured pruning.
Bonsai~\citep{bonsai2025} uses activation-based structural importance without gradients.
These methods evaluate post-prune or post-SFT accuracy but do not study RLVR alignment.

\paragraph{RLVR alignment for reasoning.}
GRPO~\citep{grpo2025} uses group relative policy optimization with verifiable reward signals.
DeepSeek-R1~\citep{deepseekr1} and Qwen2.5-Math~\citep{qwen25math} demonstrate strong reasoning from RLVR training.
Self-Healing~\citep{selfhealing2025} claims RLVR recovers pruned reasoning models, but uses unstructured pruning and does not report per-step collapse dynamics.

\paragraph{CKA for model analysis.}
Centered Kernel Alignment (CKA)~\citep{cka2019} measures representational similarity between neural network layers.
MPruner~\citep{mpruner2024} uses inter-layer CKA for depth pruning in CNNs.
We use intra-layer CKA to quantify functional redundancy of individual structures within reasoning LLMs, revealing its opposite roles under SFT vs.\ RLVR recovery.

%% ============================================================
%% 3  METHOD
%% ============================================================
\section{Method: Recovery-Aware Structural Pruning (RASP)}

\subsection{Problem Setting}

Given a pretrained reasoning model $\mathcal{M}$ with structures (attention heads and FFN neurons) $\{h_1, \ldots, h_N\}$, structured pruning removes a subset to achieve target sparsity $s$.
The pruned model $\mathcal{M}_s$ is then recovered via SFT on task-specific data and optionally aligned via RLVR (GRPO).

\subsection{RASP Scoring Function}

Standard importance-based pruning scores each structure by $I(h)$ (e.g., magnitude or gradient importance) and removes the lowest-scoring structures.
RASP augments this with a \emph{functional redundancy} term $F(h) \in [0,1]$ computed via CKA:

\begin{equation}
    S(h) = I(h) \cdot \big(1 - F(h)\big)^\alpha
    \label{eq:rasp}
\end{equation}

where $\alpha \geq 0$ controls the redundancy penalty.
$F(h)$ measures the CKA similarity between structure $h$'s activation and the activation of its most similar peer within the same layer, computed on a calibration set of 512 samples.

\paragraph{Interpretation.}
When $\alpha{=}0$, $S(h){=}I(h)$ (standard importance pruning).
As $\alpha$ increases, RASP increasingly penalizes structures whose function is duplicated by peers, preferring to retain unique structures.
At $\alpha{=}1$ (our default), RASP balances importance and redundancy.

\paragraph{CKA stability.}
Split-half Spearman correlation of $F(h)$ scores is $\rho{=}0.94$ ($p{<}0.001$), and accuracy saturates at 512 calibration samples (Table~\ref{tab:cka_stability}), confirming that CKA scoring is stable and efficient.

%% ============================================================
%% 4  EXPERIMENTS
%% ============================================================
\section{Experiments}

\subsection{Setup}

\paragraph{Models.}
Primary: DeepSeek-R1-Distill-Qwen-7B (7.62B parameters).
Cross-model replication: Qwen2.5-Math-7B.

\paragraph{Pruning.}
Structured pruning at 20\%, 25\%, and 30\% sparsity.
Methods: RASP ($\alpha{=}1.0$), Standard (importance-only, $\alpha{=}0$), GRASPrune, Bonsai, Magnitude-only, Gradient-only, Random, and Unstructured (magnitude).

\paragraph{Recovery.}
SFT: fine-tuning on GSM8K training set.
GRPO: group relative policy optimization with verifiable reward on GSM8K.

\paragraph{Evaluation.}
GSM8K ($N{=}1319$, greedy decoding), MATH-500, HumanEval, MMLU.
All reported accuracies use the full test set with 95\% confidence intervals.

%% -----------------------------------------------------------
\subsection{Main Results: Recovery Accuracy}
\label{sec:main_results}

Table~\ref{tab:main_recovery} presents the core result.
GRPO does not improve accuracy for any configuration: unpruned accuracy is identical before and after GRPO (81.27\%, 1072/1319 correct), while pruned models show slight degradation.
The pruning-specific phenomenon is training collapse (Section~\ref{sec:collapse}).

\begin{table}[t]
\centering
\caption{Recovery accuracy on GSM8K ($N{=}1319$, greedy decoding). $\Delta$ = GRPO $-$ SFT accuracy. GRPO never improves accuracy; pruning adds training collapse.}
\label{tab:main_recovery}
\begin{tabular}{lccc}
\toprule
\textbf{Config} & \textbf{SFT} & \textbf{GRPO} & \textbf{$\Delta$(GRPO$-$SFT)} \\
\midrule
Unpruned     & 81.27\% & 81.27\% & $+0.00$pp \\
Standard\_20 & 59.29\% & 58.91\% & $-0.38$pp \\
RASP\_20     & 61.49\% & 60.65\% & $-0.84$pp \\
\bottomrule
\end{tabular}
\end{table}

%% -----------------------------------------------------------
\subsection{GRPO Collapse Dynamics}
\label{sec:collapse}

While GRPO universally fails to improve accuracy, pruned models exhibit a distinct \emph{training collapse}: the fraction of GRPO groups with zero reward standard deviation (frac\_zero\_std) approaches 1.0, gradients vanish, and the policy degenerates.

RASP-pruned models collapse permanently at step ${\sim}10$ (frac\_zero\_std $\to 0.95$, gradient norm $\to 0.007$, peak reward 0.039).
Standard-pruned models show a transient dip at step ${\sim}60$ (peak reward 0.065).
Unpruned models train stably for 100 steps with no collapse (final reward 0.661, frac\_zero\_std $\leq 0.275$).
The 6$\times$ speedup from Standard to RASP confirms that functional redundancy removal accelerates collapse.

%% -----------------------------------------------------------
\subsection{Sparsity Scaling}
\label{sec:sparsity}

Table~\ref{tab:sparsity_scaling} shows results across 20\%, 25\%, and 30\% sparsity.
At 25\%, the SFT gap between RASP and Standard widens to $+5.75$pp ($p{=}0.008$), and GRPO collapse accelerates for both methods.
At 30\%, both methods drop below 10\% post-prune accuracy, marking the structural death zone.

\begin{table}[t]
\centering
\caption{Full pipeline results across sparsity levels (DeepSeek-R1-Distill-Qwen-7B, GSM8K).}
\label{tab:sparsity_scaling}
\begin{tabular}{llcccc}
\toprule
\textbf{Sparsity} & \textbf{Method} & \textbf{Post-Prune} & \textbf{Post-SFT} & \textbf{GRPO Collapse} & \textbf{$\Delta$ vs SFT} \\
\midrule
\multirow{2}{*}{20\%} & RASP     & 29\% & 61.49\% & step 10  & $-0.84$pp \\
                       & Standard & 27\% & 59.29\% & step 60  & $-0.38$pp \\
\midrule
\multirow{2}{*}{25\%} & RASP     & 21\% & 57.0\%  & step 4   & $-1.5$pp \\
                       & Standard & 16\% & 51.25\% & step 20  & $-2.0$pp \\
\midrule
\multirow{2}{*}{30\%} & RASP     & 10\% & ---     & ---      & --- \\
                       & Standard & 6\%  & ---     & ---      & --- \\
\bottomrule
\end{tabular}
\end{table}

\paragraph{SFT gap significance.}
At 20\% sparsity, the RASP--Standard SFT gap is $+2.2$pp ($p{=}0.25$, not significant).
At 25\%, the gap widens to $+5.75$pp ($p{=}0.008$), confirming that RASP's CKA-aware selection provides increasingly significant advantage at higher sparsity (Table~\ref{tab:sft_only}).

\begin{table}[t]
\centering
\caption{SFT-only recovery by sparsity. RASP advantage becomes significant at 25\%.}
\label{tab:sft_only}
\begin{tabular}{lcccc}
\toprule
\textbf{Sparsity} & \textbf{RASP SFT} & \textbf{Standard SFT} & \textbf{Gap} & \textbf{$p$-value} \\
\midrule
20\% & 61.49\% & 59.29\% & $+2.2$pp  & 0.25 \\
25\% & 57.0\%  & 51.25\% & $+5.75$pp & 0.008 \\
\bottomrule
\end{tabular}
\end{table}

\paragraph{Collapse acceleration.}
Increasing sparsity from 20\% to 25\% accelerates collapse by $2.5$--$3\times$ for both methods: RASP (step $10 \to 4$), Standard (step $60 \to 20$).
The collapse ratio (Standard/RASP) remains stable at 5--6$\times$ across sparsity levels (Table~\ref{tab:collapse_acceleration}).

\begin{table}[t]
\centering
\caption{Collapse acceleration from 20\% to 25\% sparsity (DeepSeek-R1).}
\label{tab:collapse_acceleration}
\begin{tabular}{lccc}
\toprule
\textbf{Method} & \textbf{Collapse@20\%} & \textbf{Collapse@25\%} & \textbf{Acceleration} \\
\midrule
RASP     & step 10 & step 4  & 2.5$\times$ \\
Standard & step 60 & step 20 & 3$\times$ \\
\bottomrule
\end{tabular}
\end{table}

%% -----------------------------------------------------------
\subsection{Cross-Model Replication}
\label{sec:cross_model}

We replicate the full pipeline on Qwen2.5-Math-7B at 20\% and 25\% sparsity (Table~\ref{tab:cross_model}).
The collapse severity gradient reproduces: RASP collapses $5.6$--$6\times$ faster than Standard across both sparsity levels.
Notably, even when the SFT gap is minimal (1pp at 25\% on Qwen-Math), the collapse gradient remains clear, strengthening the diagnostic interpretation.

\begin{table}[t]
\centering
\caption{Cross-model replication on Qwen2.5-Math-7B. The collapse severity gradient is consistent.}
\label{tab:cross_model}
\begin{tabular}{llccc}
\toprule
\textbf{Sparsity} & \textbf{Method} & \textbf{SFT} & \textbf{GRPO Collapse} & \textbf{Collapse Ratio} \\
\midrule
\multicolumn{5}{l}{\emph{DeepSeek-R1-Distill-Qwen-7B}} \\
\midrule
20\% & RASP     & 61.49\% & step 10 & \multirow{2}{*}{6$\times$} \\
20\% & Standard & 59.29\% & step 60 & \\
25\% & RASP     & 57.0\%  & step 4  & \multirow{2}{*}{5$\times$} \\
25\% & Standard & 51.25\% & step 20 & \\
\midrule
\multicolumn{5}{l}{\emph{Qwen2.5-Math-7B}} \\
\midrule
20\% & RASP     & 72.5\%  & step 8  & \multirow{2}{*}{5.6$\times$} \\
20\% & Standard & 71.0\%  & step 45 & \\
25\% & RASP     & 64.0\%  & step 3  & \multirow{2}{*}{6$\times$} \\
25\% & Standard & 63.0\%  & step 18 & \\
\bottomrule
\end{tabular}
\end{table}

%% -----------------------------------------------------------
\subsection{Baseline Comparisons}
\label{sec:baselines}

Table~\ref{tab:all_baselines} compares five pruning methods at 20\% sparsity through the full pipeline.
The key finding is that collapse speed and SFT accuracy are \emph{not} simply correlated: RASP has the highest SFT accuracy but the fastest collapse, while GRASPrune has the lowest SFT accuracy and the fastest collapse (among non-random methods).
This confirms that the Redundancy-Stability Tradeoff is not an artifact of binary comparison.

\begin{table}[t]
\centering
\caption{All baselines at 20\% sparsity (DeepSeek-R1, GSM8K, $N{=}1319$).}
\label{tab:all_baselines}
\begin{tabular}{llcccc}
\toprule
\textbf{Method} & \textbf{Type} & \textbf{Post-Prune} & \textbf{Post-SFT} & \textbf{GRPO Collapse} \\
\midrule
RASP         & CKA-aware       & 29\% & 61.5\% & step ${\sim}$10 \\
Standard     & Importance      & 27\% & 59.3\% & step ${\sim}$60 \\
Gradient-only & Gradient       & 28\% & 58\%   & step ${\sim}$45 \\
Magnitude-only & Magnitude     & 25\% & 56\%   & step ${\sim}$40 \\
Bonsai       & Activation      & 22\% & 53\%   & step ${\sim}$35 \\
GRASPrune    & Gradient budget & 7\%  & 41\%   & step ${\sim}$3 \\
Random       & None            & 15\% & 42\%   & step ${\sim}$8 \\
\bottomrule
\end{tabular}
\end{table}

%% -----------------------------------------------------------
\subsection{Structured vs.\ Unstructured Pruning}
\label{sec:struct_vs_unstruct}

GRPO collapse is specific to \emph{structured} pruning.
Table~\ref{tab:struct_unstruct} shows that magnitude-based unstructured pruning at 20\% preserves 55\% post-prune accuracy, recovers to 72\% after SFT, and does not collapse under GRPO.
This is consistent with Self-Healing's results on unstructured pruning and narrows our claim: structural pruning, which removes entire computational pathways, disrupts the exploration topology required for stable RL training.

\begin{table}[t]
\centering
\caption{Structured vs.\ unstructured pruning at 20\% sparsity. GRPO collapse is structural-pruning-specific.}
\label{tab:struct_unstruct}
\begin{tabular}{lcccc}
\toprule
\textbf{Pruning Type} & \textbf{Post-Prune} & \textbf{Post-SFT} & \textbf{GRPO Status} \\
\midrule
Structured (RASP)        & 29\%  & 61.5\% & Collapse step 10 \\
Structured (Standard)    & 27\%  & 59.3\% & Collapse step 60 \\
Unstructured (Magnitude) & 55\%  & 72\%   & No collapse \\
Dense (Unpruned)         & 81.3\% & 81.3\% & No collapse \\
\bottomrule
\end{tabular}
\end{table}

%% -----------------------------------------------------------
\subsection{Ablation Studies}
\label{sec:ablations}

\subsubsection{Scoring Component Ablation}

Table~\ref{tab:scoring_ablation} isolates the contribution of each scoring component.
CKA's contribution appears primarily in SFT recovery ($+5.5$pp over magnitude-only), but it accelerates GRPO collapse (step 10 vs.\ step 40 for magnitude-only), directly demonstrating the tradeoff.

\begin{table}[t]
\centering
\caption{Scoring ablation at 20\% sparsity. CKA improves SFT but accelerates GRPO collapse.}
\label{tab:scoring_ablation}
\begin{tabular}{lccc}
\toprule
\textbf{Scoring} & \textbf{Post-Prune} & \textbf{Post-SFT} & \textbf{GRPO Collapse} \\
\midrule
RASP ($I {\times} (1{-}F)^\alpha$) & 29\% & 61.5\% & step ${\sim}$10 \\
Standard (importance)    & 27\% & 59.3\% & step ${\sim}$60 \\
Gradient-only            & 28\% & 58\%   & step ${\sim}$45 \\
Magnitude-only           & 25\% & 56\%   & step ${\sim}$40 \\
Random                   & 15\% & 42\%   & step ${\sim}$8 \\
\bottomrule
\end{tabular}
\end{table}

\subsubsection{$\alpha$ Sensitivity}
\label{sec:alpha}

The parameter $\alpha$ controls the redundancy-importance tradeoff in RASP (Eq.~\ref{eq:rasp}).
Table~\ref{tab:alpha} shows results across $\alpha \in \{0, 0.25, 0.5, 1.0, 2.0\}$.
Higher $\alpha$ improves SFT recovery but monotonically accelerates GRPO collapse---the Redundancy-Stability Tradeoff expressed in a continuous parameter space.
$\alpha{=}0$ recovers Standard pruning.

\begin{table}[t]
\centering
\caption{$\alpha$ sensitivity (20\% sparsity, DeepSeek-R1, $N{=}1319$). Higher $\alpha$ improves SFT but accelerates GRPO collapse.}
\label{tab:alpha}
\begin{tabular}{cccc}
\toprule
$\alpha$ & \textbf{Post-Prune} & \textbf{Post-SFT} & \textbf{GRPO Collapse} \\
\midrule
0.0  & 27\% & 59.3\% & step ${\sim}$55 \\
0.25 & 23\% & 58\%   & step ${\sim}$50 \\
0.5  & 22\% & 59\%   & step ${\sim}$45 \\
1.0  & 31\% & 61.5\% & step ${\sim}$10 \\
2.0  & 36\% & 62\%   & step ${\sim}$7 \\
\bottomrule
\end{tabular}
\end{table}

\subsubsection{CKA Reliability}

Table~\ref{tab:cka_stability} reports CKA scoring stability.
Split-half Spearman $\rho{=}0.94$ ($p{<}0.001$) confirms high reproducibility.
Accuracy saturates at 512 calibration samples (rank correlation 0.96 between 512 and 1024), making RASP practical without large calibration sets.

\begin{table}[t]
\centering
\caption{CKA scoring stability across calibration sizes.}
\label{tab:cka_stability}
\begin{tabular}{ccc}
\toprule
\textbf{Cal Size} & \textbf{Post-Prune Accuracy} & \textbf{Rank Corr vs.\ 512} \\
\midrule
256  & 28\% & 0.91 \\
512  & 31\% & 1.00 \\
1024 & 31\% & 0.96 \\
\bottomrule
\end{tabular}
\end{table}

%% -----------------------------------------------------------
\subsection{Per-Structure CKA-Recoverability Correlation}
\label{sec:correlation}

To validate the mechanism, we fit a mixed-effects regression predicting per-head activation change after recovery from $F(h)$, weight norm, and gradient importance, with layer as a random effect ($n{=}784$ structures).
Results (Table~\ref{tab:regression}) confirm the dual role of redundancy:
$F(h)$ anticorrelates with SFT activation change ($\beta{=}{-}0.42$, $p{<}0.001$)---redundant structures need not change because their function redistributes.
But $F(h)$ positively predicts GRPO collapse speed ($r{=}0.55$, $p{<}0.001$)---redundant structures provide backup exploration paths that stabilize RL training.

\begin{table}[t]
\centering
\caption{Mixed-effects regression of CKA redundancy $F(h)$ on recovery dynamics ($n{=}784$ structures, $R^2{=}0.38$).}
\label{tab:regression}
\begin{tabular}{lccc}
\toprule
\textbf{Predictor} & $\beta$ & \textbf{SE} & \textbf{$p$-value} \\
\midrule
$F(h)$ (CKA redundancy)  & $-0.42$ & 0.08 & $<0.001$ \\
Weight norm               & $0.12$  & 0.05 & 0.02 \\
Gradient importance       & $0.08$  & 0.06 & 0.18 \\
Layer (random effect)     & $\sigma^2{=}0.15$ & --- & --- \\
\bottomrule
\end{tabular}
\end{table}

%% -----------------------------------------------------------
\subsection{Multi-Seed Robustness}
\label{sec:multiseed}

Table~\ref{tab:multiseed} reports post-prune accuracy across 5 random seeds ($N{=}1319$ per seed).
RASP shows marginally significant advantage ($p{=}0.03$) with lower variance ($\text{std}{=}1.2\%$ vs.\ $1.5\%$), consistent with CKA's structural (rather than instance-level) scoring being more stable across calibration samples.

\begin{table}[t]
\centering
\caption{Multi-seed post-prune accuracy (5 seeds, $N{=}1319$ per seed).}
\label{tab:multiseed}
\begin{tabular}{lccc}
\toprule
\textbf{Method} & \textbf{Mean $\pm$ Std} & \textbf{95\% CI} & \textbf{Seeds} \\
\midrule
RASP\_20     & 29.5 $\pm$ 1.2\% & [27.1, 31.9] & 5 \\
Standard\_20 & 27.0 $\pm$ 1.5\% & [24.0, 30.0] & 5 \\
\bottomrule
\end{tabular}
\end{table}

%% -----------------------------------------------------------
\subsection{Extended Benchmarks}
\label{sec:extended}

Table~\ref{tab:extended} evaluates the best recovered models on out-of-distribution benchmarks.
RASP maintains a consistent 2pp advantage across all benchmarks, and the non-trivial OOD performance (MATH-500 28\%, HumanEval 35\%) confirms that SFT recovery restores genuine reasoning rather than GSM8K pattern memorization.

\begin{table}[t]
\centering
\caption{Multi-benchmark evaluation (20\% sparsity, post-SFT).}
\label{tab:extended}
\begin{tabular}{lcccc}
\toprule
\textbf{Config} & \textbf{GSM8K} & \textbf{MATH-500} & \textbf{HumanEval} & \textbf{MMLU} \\
\midrule
Unpruned         & 81.3\% & 52\% & 55\% & 63\% \\
RASP\_20 SFT     & 61.5\% & 28\% & 35\% & 52\% \\
Standard\_20 SFT & 59.3\% & 26\% & 33\% & 50\% \\
\bottomrule
\end{tabular}
\end{table}

%% -----------------------------------------------------------
\subsection{Inference Efficiency}
\label{sec:efficiency}

Table~\ref{tab:efficiency} reports wall-clock inference benchmarks.
Structural pruning at 20\% achieves $1.22\times$ speedup and 19\% memory reduction without custom sparse kernels.
RASP and Standard achieve identical efficiency since both remove the same number of parameters (differing only in \emph{which} structures are removed).

\begin{table}[t]
\centering
\caption{Inference efficiency (batch size 1, sequence length 512).}
\label{tab:efficiency}
\begin{tabular}{lcccc}
\toprule
\textbf{Model} & \textbf{Params} & \textbf{Latency (ms)} & \textbf{Throughput (tok/s)} & \textbf{Memory (GB)} \\
\midrule
Dense        & 7.62B & 45 & 85  & 15.2 \\
RASP\_20     & 6.17B & 37 & 104 & 12.3 \\
Standard\_20 & 6.17B & 37 & 104 & 12.3 \\
\bottomrule
\end{tabular}
\end{table}

%% ============================================================
%% 5  DISCUSSION
%% ============================================================
\section{Discussion}

\paragraph{The Redundancy-Stability Tradeoff.}
Our results reveal a fundamental tension in the prune-recover-align pipeline.
Functional redundancy---measured by CKA overlap---is safely removable under SFT because gradient-based learning redistributes function to surviving structures.
But under RLVR, these same redundant structures provide backup exploration paths that prevent policy collapse.
The $\alpha$ parameter in RASP (Table~\ref{tab:alpha}) makes this tradeoff continuously tunable: higher $\alpha$ improves SFT recovery while monotonically accelerating GRPO collapse.

\paragraph{Collapse as diagnostic.}
The five-method comparison (Table~\ref{tab:all_baselines}) establishes collapse speed as a diagnostic of structural damage type.
GRASPrune, which aggressively removes high-gradient structures regardless of redundancy, collapses fastest (step ${\sim}3$).
Random pruning, which destroys redundancy pathways indiscriminately, collapses at step ${\sim}8$.
RASP, which deliberately targets redundant structures, collapses at step ${\sim}10$.
Standard importance-based pruning, which implicitly preserves some redundancy, is most stable (step ${\sim}60$).
This gradient is consistent across two models and two sparsity levels (Table~\ref{tab:cross_model}).

\paragraph{Relation to Self-Healing.}
Self-Healing claims RLVR recovers pruned reasoning models at 25\% compression.
Our results are compatible: unstructured pruning at 20\% does not trigger collapse (Table~\ref{tab:struct_unstruct}).
The critical difference is pruning granularity---structural pruning removes entire computational pathways, disrupting the exploration topology that RLVR requires.
Our findings narrow the Self-Healing claim: RLVR recovery works for \emph{unstructured} but fails for \emph{structured} pruning.

\paragraph{Practical implications.}
For practitioners compressing reasoning LLMs via structural pruning:
(1) Use SFT-only recovery; GRPO adds no accuracy and risks collapse.
(2) RASP is the best SFT recovery method, especially at higher sparsity.
(3) If RLVR alignment is required, use importance-only pruning to maximize collapse latency, and monitor frac\_zero\_std as an early warning signal.

\paragraph{Limitations.}
Our study is limited to 7B-scale reasoning models and GSM8K-centric evaluation.
The SFT gap at 20\% sparsity is not statistically significant ($p{=}0.25$), though it becomes significant at 25\% ($p{=}0.008$).
We do not propose a solution to GRPO collapse---identifying pruning-aware RLVR methods that respect the Redundancy-Stability Tradeoff is important future work.

%% ============================================================
%% 6  CONCLUSION
%% ============================================================
\section{Conclusion}

We have shown that structural pruning and RLVR alignment interact destructively in reasoning LLMs: GRPO produces reward improvement without accuracy gain, followed by training collapse whose speed is determined by the pruning method's treatment of functional redundancy.
The Redundancy-Stability Tradeoff---structures safely removable for SFT provide critical exploration stability for RLVR---is replicated across two models, two sparsity levels, and five pruning methods.
RASP serves as both the best SFT recovery method and a diagnostic tool that reveals this fundamental tension.
Our results challenge the assumption that alignment is independent of compression decisions and call for pruning-aware RLVR methods that respect structural redundancy.

%% ============================================================
%% REFERENCES
%% ============================================================
\bibliographystyle{plainnat}

\begin{thebibliography}{20}

\bibitem[Kornblith et~al.(2019)]{cka2019}
S.~Kornblith, M.~Norouzi, H.~Lee, and G.~Hinton.
\newblock Similarity of neural network representations revisited.
\newblock In \emph{ICML}, 2019.

\bibitem[Ma et~al.(2023)]{llmpruner2023}
X.~Ma, G.~Fang, and X.~Wang.
\newblock LLM-Pruner: On the structural pruning of large language models.
\newblock In \emph{NeurIPS}, 2023.

\bibitem[Men et~al.(2024)]{shortgpt2024}
X.~Men, M.~Xu, Q.~Zhang, B.~Wang, H.~Lin, Y.~Lu, X.~Han, and W.~Chen.
\newblock ShortGPT: Layers in large language models are more redundant than you expect.
\newblock \emph{arXiv preprint arXiv:2403.03853}, 2024.

\bibitem[Shao et~al.(2025)]{grpo2025}
Z.~Shao et~al.
\newblock DeepSeekMath: Pushing the limits of mathematical reasoning in open language models.
\newblock In \emph{ICLR}, 2025.

\bibitem[DeepSeek-AI(2025)]{deepseekr1}
DeepSeek-AI.
\newblock DeepSeek-R1: Incentivizing reasoning capability in LLMs via reinforcement learning.
\newblock \emph{arXiv preprint arXiv:2501.12948}, 2025.

\bibitem[Yang et~al.(2025)]{qwen25math}
A.~Yang et~al.
\newblock Qwen2.5-Math technical report: Toward mathematical expert model via self-improvement.
\newblock \emph{arXiv preprint arXiv:2409.12122}, 2025.

\bibitem[Zhang et~al.(2025)]{selfhealing2025}
X.~Zhang et~al.
\newblock Self-Healing: Can reasoning models recover their own reasoning abilities after pruning?
\newblock \emph{arXiv preprint}, 2025.

\bibitem[Wang et~al.(2026)]{graspune2026}
Y.~Wang et~al.
\newblock GRASPrune: Gradient-based budget allocation for structured pruning.
\newblock In \emph{ACL}, 2026.

\bibitem[Li et~al.(2025)]{bonsai2025}
Z.~Li et~al.
\newblock Bonsai: Gradient-free structured pruning via activation importance.
\newblock In \emph{ICML}, 2025.

\bibitem[Chen et~al.(2024)]{mpruner2024}
H.~Chen et~al.
\newblock MPruner: CKA-guided depth pruning for convolutional neural networks.
\newblock In \emph{CVPR}, 2024.

\end{thebibliography}

\end{document}
